% ---
% id: transmitter/rem:certificate-forms
% kind: rem
% title: "Purity and exact declarability of the certificate forms"
% summary: >-
%   Only the converse in part (ii) of [[thm:certificate-sufficiency]] depends on purity. A
%   certificate vector, a globally pure operator over sectors and a pure block
%   over non-orthogonal states can all need irrational entries, which is why
%   the certificate is an operator with a block form and an ensemble form.
% prominence: standard
% status: accepted
% uses: [def:srcert]
% mentions: [thm:certificate-sufficiency, cor:noncanonical, lem:ensemble-presentation, def:source-spec]
% assumes: []
% refs: []
% provenance: newly-derived
% lean:
% ---

\begin{remark}\label{rem:certificate-forms}
\emph{What purity adds.} A pure-form certificate satisfies every condition of
the block form at the one-sector structure and one condition more
(Definition~\ref{def:srcert}). Every statement this module proves of the
block form therefore holds of the pure form, and
Theorem~\ref{thm:certificate-sufficiency}\,(ii) is the only statement whose \emph{conclusion}
depends on the additional condition: in pure form a source/POVM mismatch of
the kind (ii) describes is a detected defect and not a silent one, while in
block form it is the POVM condition of Definition~\ref{def:source-spec},
decided on its own, that excludes it. Part (i) of the theorem lists purity as
a condition of the pure form, and a consumer that demands purity of every
certificate re-decides it and therefore refuses a block form. Everything
Corollary~\ref{cor:noncanonical} exhibits is admitted in block form as well.
Not every block-form certificate is pure: already at one sector the block
form admits a positive semidefinite operator of unit trace and rank above
one, which the pure form refuses.

\emph{Why the certificate is an operator and not a vector.} No binary64
number $a$ satisfies $a^2=\tfrac12$, so a declared vector with entries
$\mathrm{fl}(1/\sqrt2)$ yields $\dyad{\psi}$ with entries
$\mathrm{fl}(1/\sqrt2)^2\ne\tfrac12$, and (SR1) fails exactly \emph{for the
certificate as an instance declares it} --- $\ket{\Phip}$ in that
instance's own computational basis, whose amplitudes are $1/\sqrt2$. The
\emph{operator} $\dyad{\Phip}$, by contrast, has every entry equal to
$\tfrac12$, which is a binary64 number; and every quantity (SR1), (SR2)
and Theorem~\ref{thm:certificate-sufficiency} read off the certificate is a function of
$\Psi$ and never of a vector. In pure form the purity conditions ---
$\Psi^2=\Psi$ with $\tr\Psi=1$ --- are what recover the vector where an
argument needs one (Theorem~\ref{thm:certificate-sufficiency}\,(ii)), and they are two exact
identities on the declared matrix. Whether a given declaration of the
canonical data of Corollary~\ref{cor:noncanonical} is admissible is decided,
not assumed: the canonical $\Psi$ must be declared as a binary64 rank-one
projector, and a declaration whose entries do not satisfy $\Psi^2=\Psi$
exactly is refused.

\emph{Why a block form and an ensemble presentation.} Exact purity is not
declarable for a source that draws a sector and then prepares a state in it.
A globally pure $\Psi$ over sectors of weights $g_t$ carries cross-sector
amplitudes $\sqrt{g_tg_{t'}}$, which need not be dyadic rationals, and a
pure block whose marginal has the spectrum $(\tfrac12,\tfrac14,\tfrac14)$
carries the entry $\sqrt{1/8}$, which is not one. The direct sum of positive
blocks is exactly declarable, it is the state of the source that first draws
$t$ with probability $\tr\Psi_t$ and then prepares $\Psi_t/\tr\Psi_t$, and
the sufficiency statement, Theorem~\ref{thm:certificate-sufficiency}\,(iii), reads purity
nowhere. Where even the blocks are not declarable --- a pure block over
non-orthogonal signal states carries square roots of their weights --- the
ensemble presentation is the form to use: it has the purity
(Lemma~\ref{lem:ensemble-presentation}\,(ii)) and declares none of the
numbers, which is what it exists for.
\end{remark}
